% Source: https://artofproblemsolving.com/community/c3h1873632_base_question
% Stable identity: aops-topic-1873632-q01
\textbf{(a)} By definition, the base-$4$ numeral $2123_4$ stands for the number $2\cdot4^3+1\cdot4^2+2\cdot4^1+3\cdot4^0$: the digits count units, fours, sixteens and sixty-fours. This is a statement about \emph{numbers}, not about any way of writing them. We do the arithmetic in base ten, so the result comes out as a base-ten numeral:
\[
2\cdot64+1\cdot16+2\cdot4+3=128+16+8+3=155 .
\]
\textbf{(b)} The powers of $4$ are $1,4,16,64,256,1024$. Repeatedly take the largest multiples:
\[
2123=2\cdot1024+75,\quad 75=0\cdot256+1\cdot64+11,\quad 11=0\cdot16+2\cdot4+3,
\]
so $2123=2\cdot4^5+0\cdot4^4+1\cdot4^3+0\cdot4^2+2\cdot4+3=201023_4$.

\textbf{(c)} The numeral $2123$ (base ten) means $2\cdot\text{ten}^3+1\cdot\text{ten}^2+2\cdot\text{ten}+3$, and ten is $22_4$. Working in base $4$:
\[
(22_4)^2=1210_4,\qquad(22_4)^3=1210_4\cdot22_4=33220_4,
\]
\[
2\cdot33220_4=133100_4,\qquad 1\cdot1210_4=1210_4,\qquad 2\cdot22_4=110_4,
\]
and adding in base $4$: $133100_4+1210_4=200310_4$, then $200310_4+110_4=201020_4$, then $201020_4+3_4=201023_4$. This is the answer to (b). So the poster's idea is right: the same expansion converts \emph{into} whichever base the arithmetic is done in.

Answers: $\boxed{2123_4=155,\quad2123=201023_4}$.
